\documentclass[11pt,a4paper]{article}
\usepackage[T1]{fontenc}
\usepackage[utf8]{inputenc}
\usepackage{amsmath,amssymb,xcolor}
\usepackage[margin=25mm]{geometry}
\usepackage[colorlinks=true,urlcolor=blue]{hyperref}
\definecolor{nexorablue}{HTML}{2459B5}
\setlength{\parindent}{0pt}
\setlength{\parskip}{7pt}
\setlength{\emergencystretch}{2em}
\renewcommand{\familydefault}{\sfdefault}
\begin{document}

{\small\bfseries\color{nexorablue} NEXORA CPGE}\par
{\small MPSI / Cours}\par
{\LARGE\bfseries Lois générales des circuits électriques}\par
\vspace{0.25em}\hrule\vspace{0.6em}

\section*{Quand peut-on négliger la propagation ?}
Une variation électrique met un temps de l’ordre de $\tau_p=\ell/v$ à parcourir un circuit de taille caractéristique $\ell$, avec $v$ la vitesse de propagation dans le milieu. Dans l’approximation des régimes quasi-stationnaires (ARQS), ce temps est très petit devant la durée caractéristique de variation. Pour un signal périodique de fréquence $f$, on compare $\ell f/v$ à $1$. Un front bref impose de comparer $\tau_p$ à la durée du front, même si les impulsions se répètent lentement.

Exemple. Pour $\ell=30\,\mathrm{cm}$ et $v=2{,}0\times10^8\,\mathrm{m\,s^{-1}}$, on obtient $\tau_p=1{,}5\,\mathrm{ns}$. Ce temps est très inférieur à la période d’un signal sinusoïdal de $10\,\mathrm{kHz}$. On peut utiliser un même courant instantané le long d’une branche sans dérivation dans le modèle usuel.

\section*{Charge, débit et signe du courant}
La charge $q$ se mesure en coulombs. On choisit un sens de référence à travers une section du conducteur. L’intensité algébrique $i$ est le débit de charge dans ce sens, soit $i=\mathrm dq_{\mathrm{tr}}/\mathrm dt$ si $q_{\mathrm{tr}}$ désigne la charge algébrique ayant traversé la section. Un ampère vaut un coulomb par seconde. Le courant conventionnel suit le transport de charges positives et s’oppose au mouvement des électrons dans un métal.

Sur un intervalle de temps, $\Delta q_{\mathrm{tr}}=\int i(t)\,\mathrm dt$. Un courant constant de $8\,\mathrm{mA}$ pendant $0{,}25\,\mathrm s$ transporte une charge algébrique de $2\,\mathrm{mC}$ dans le sens choisi. Un courant négatif se lit dans le sens inverse de sa flèche.

\section*{Nœuds et conservation de la charge}
Un nœud modélise une jonction entre branches. Pour un petit volume entourant la jonction et contenant la charge $q_N$, le taux de variation $\mathrm dq_N/\mathrm dt$ est la différence entre les courants entrants et sortants. Dans le modèle de nœud idéal en ARQS, l’accumulation est négligeable, d’où
\[ \sum_{\mathrm{entrants}} i_k=\sum_{\mathrm{sortants}} i_k. \]
Un courant affecté à un côté de l’équation conserve son signe algébrique. Exemple. Deux courants orientés vers un nœud valent $7$ et $-2\,\mathrm{mA}$. Un courant orienté vers l’extérieur vaut alors $5\,\mathrm{mA}$ si ces trois branches sont les seules.

\section*{Potentiel, masse et tension}
On choisit un point de référence auquel on attribue le potentiel $0$. Le potentiel d’un autre point dépend de cette référence. La tension $U_{AB}=V_A-V_B$ reste la même si tous les potentiels sont déplacés d’une même constante. La flèche de $U_{AB}$ pointe de $B$ vers $A$ dans la convention graphique utilisée ici.

Exemple. Avec $V_A=-1\,\mathrm V$ et $V_B=3\,\mathrm V$, on a $U_{AB}=-4\,\mathrm V$ et $U_{BA}=4\,\mathrm V$. Des points reliés par un fil idéal ont le même potentiel. La masse choisie pour le calcul est une référence de potentiel, et sa liaison éventuelle à la terre dépend du montage.

\section*{Maille et différences de potentiel}
Dans le modèle de circuit à dipôles, les différences de potentiel le long d’un parcours fermé se compensent. Pour la boucle $A,B,C,A$,
\[ U_{AB}+U_{BC}+U_{CA}=0. \]
La flèche d’une tension et le sens de parcours sont deux choix distincts. Pour éviter les erreurs, écrire chaque terme avec les lettres dans l’ordre du parcours, puis traduire via la loi du dipôle. Exemple. Si $U_{AB}=6\,\mathrm V$ et $U_{BC}=-2\,\mathrm V$, alors $U_{AC}=4\,\mathrm V$ et $U_{CA}=-4\,\mathrm V$.

\section*{Conventions et caractère d’un dipôle}
Soit $u=V_A-V_B$. En convention récepteur, le courant est orienté de $A$ vers $B$ dans le dipôle. Alors $p_{\text{reçue}}=ui$. En convention générateur, il est orienté de $B$ vers $A$ et $p_{\text{reçue}}=-ui$. Dans cette seconde convention, $ui$ est la puissance fournie algébrique.

Le caractère physique se déduit du signe de la puissance reçue, une fois la convention appliquée. Si $p_{\text{reçue}}>0$, le dipôle reçoit de l’énergie électrique. Si $p_{\text{reçue}}<0$, il en cède. Une batterie peut passer d’un rôle à l’autre lors de la décharge et de la recharge.

\section*{Puissance instantanée et énergie}
Un volt-ampère vaut un watt pour le produit instantané $u(t)i(t)$. L’énergie reçue entre $t_1$ et $t_2$ vaut $W=\int_{t_1}^{t_2}p_{\text{reçue}}(t)\,\mathrm dt$, en joules. Lorsque la puissance est constante, on multiplie par la durée exprimée en secondes.

Dans un circuit idéal fermé, la somme des puissances reçues par tous les dipôles est nulle. Pour le voir, écrire chaque contribution sous la forme $(V_A-V_B)i_{A\to B}$. Une fois les coefficients regroupés par nœud, chaque coefficient est une somme algébrique de courants qui s’annule par conservation de la charge. Ce bilan permet de vérifier un calcul de signe ou une valeur de courant.

\section*{Premières applications avec un résistor}
Le résistor idéal de résistance $R>0$ vérifie $u=Ri$ en convention récepteur. Sa puissance reçue est $Ri^2=u^2/R$. Si le courant calculé est négatif, il circule dans le sens opposé au choix initial et la dissipation reste positive. L’ohm vaut un volt par ampère.

Une branche alimentée par une source idéale de $5\,\mathrm V$ et contenant une résistance de $1\,\mathrm{k\Omega}$ est parcourue par $5\,\mathrm{mA}$ dans le sens de la chute de potentiel. La source fournit $25\,\mathrm{mW}$ et la résistance les dissipe. Sur table, vérifier l’ordre de grandeur des intensités et les limites de puissance des modèles employés.

\section*{Série, dérivation et mesure de tension}
Deux dipôles sont en série si aucune autre branche ne part de leur jonction. Ils sont alors traversés par un même courant à un signe près. Deux dipôles sont en parallèle s’ils sont reliés aux mêmes deux nœuds. Leurs tensions sont alors égales avec la même orientation. Dans un montage avec dérivation, écrire la loi du nœud avant de supposer que deux courants sont égaux.

Un voltmètre se branche entre les deux points visés. Dans le modèle idéal, il ne prélève aucun courant. Un appareil réel dont la résistance d’entrée est finie ajoute une branche et peut déplacer les potentiels. Une mesure à vide et une mesure en charge peuvent donc donner des résultats différents.

\section*{Exemple à deux branches}
Une source maintient $V_A=6\,\mathrm V$ et $V_B=0$. Deux résistances de $2\,\mathrm{k\Omega}$ et $3\,\mathrm{k\Omega}$ relient chacune $A$ à $B$. La loi de comportement fournit $3$ et $2\,\mathrm{mA}$ dans le sens $A$ vers $B$. Pour conserver la charge, la source renvoie $5\,\mathrm{mA}$ vers $A$.

Les puissances reçues par les résistances valent $18$ et $12\,\mathrm{mW}$. Comme son courant entre en $B$, la source reçoit $(V_B-V_A)\times5\,\mathrm{mA}=-30\,\mathrm{mW}$. Le bilan est nul. La même procédure s’applique quand un des courants sort avec un signe négatif.

\section*{Méthode pour un montage à plusieurs nœuds}
Nommer les nœuds, fixer la référence de potentiel et attribuer des orientations aux courants. Si la source fournit directement un potentiel par rapport à la référence, remplacer cette inconnue par sa valeur. Pour une résistance entre $A$ et $B$, le courant orienté de $A$ vers $B$ vaut $(V_A-V_B)/R$. Écrire les lois des nœuds avec les inconnues restantes, ou choisir des mailles utiles si les inconnues retenues sont des courants.

Résoudre en gardant les unités. Retrouver les courants et tensions demandés, interpréter les signes, puis vérifier une relation que l’on n’a pas utilisée dans la résolution. L’analyse d’un circuit simple ne demande pas de mémoriser une formule générale de réseau.

\section*{Fiche de révision}
ARQS : comparer le temps de propagation à la durée des variations. Au nœud : conserver la charge. Dans une maille : fermer la somme algébrique des tensions. En convention récepteur : $p_{\text{reçue}}=ui$. Pour une résistance : $p=Ri^2$. En fin d’exercice, contrôler les signes, les unités et la somme des puissances.


\end{document}
